% problem_id: S046-thm:suspbool
% source_id: S046 (HoTT Book)
% source_file: hits.tex
% statement_locator: hits.tex lines 588--591
% proof_locator: hits.tex lines 592--634
% env: lem
% pairing: tight (verified)
% status: complete (statement + author proof verbatim + reason + locators)
%
\begin{lem}\label{thm:suspbool}
  \index{type!circle}%
  $\eqv{\susp\bool}{\Sn^1}$.
\end{lem}

\begin{proof}
  Define $f:\susp\bool\to\Sn^1$ by recursion such that $f(\north)\defeq \base$ and $f(\south)\defeq\base$, while $\ap f{\merid(\bfalse)}\defid\lloop$ but $\ap f{\merid(\btrue)} \defid \refl{\base}$.
  Define $g:\Sn^1\to\susp\bool$ by recursion such that $g(\base)\defeq \north$ and $\ap g \lloop \defid \merid(\bfalse) \ct \opp{\merid(\btrue)}$.
  We now show that $f$ and $g$ are quasi-inverses.

  First we show by induction that $g(f(x))=x$ for all $x:\susp \bool$.
  If $x\jdeq\north$, then $g(f(\north)) \jdeq g(\base)\jdeq \north$, so we have $\refl{\north} : g(f(\north))=\north$.
  If $x\jdeq\south$, then $g(f(\south)) \jdeq g(\base)\jdeq \north$, and we choose the equality $\merid(\btrue) : g(f(\south)) = \south$.
  It remains to show that for any $y:\bool$, these equalities are preserved as $x$ varies along $\merid(y)$, which is to say that when $\refl{\north}$ is transported along $\merid(y)$ it yields $\merid(\btrue)$.
  By transport in path spaces and pulled back fibrations, this means we are to show that
  \[ \opp{\ap g {\ap f {\merid(y)}}} \ct \refl{\north} \ct \merid(y) = \merid(\btrue). \]
  Of course, we may cancel $\refl{\north}$.
  Now by \bool-induction, we may assume either $y\jdeq \bfalse$ or $y\jdeq \btrue$.
  If $y\jdeq \bfalse$, then we have
  \begin{align*}
    \opp{\ap g {\ap f {\merid(\bfalse)}}} \ct \merid(\bfalse)
    &= \opp{\ap g {\lloop}} \ct \merid(\bfalse)\\
    &= \opp{(\merid(\bfalse) \ct \opp{\merid(\btrue)})} \ct \merid(\bfalse)\\
    &= \merid(\btrue) \ct \opp{\merid(\bfalse)} \ct \merid(\bfalse)\\
    &= \merid(\btrue)
  \end{align*}
  while if $y\jdeq \btrue$, then we have
  \begin{align*}
    \opp{\ap g {\ap f {\merid(\btrue)}}} \ct \merid(\btrue)
    &= \opp{\ap g {\refl{\base}}} \ct \merid(\btrue)\\
    &= \opp{\refl{\north}} \ct \merid(\btrue)\\
    &= \merid(\btrue).
  \end{align*}
  Thus, for all $x:\susp \bool$, we have $g(f(x))=x$.

  Now we show by induction that $f(g(x))=x$ for all $x:\Sn^1$.
  If $x\jdeq \base$, then $f(g(\base))\jdeq f(\north)\jdeq\base$, so we have $\refl{\base} : f(g(\base))=\base$.
  It remains to show that this equality is preserved as $x$ varies along $\lloop$, which is to say that it is transported along $\lloop$ to itself.
  Again, by transport in path spaces and pulled back fibrations, this means to show that
  \[ \opp{\ap f {\ap g {\lloop}}} \ct \refl{\base} \ct \lloop = \refl{\base}.\]
  However, we have
  \begin{align*}
    \ap f {\ap g {\lloop}} &= \ap f {\merid(\bfalse) \ct \opp{\merid(\btrue)}}\\
    &= \ap f {\merid(\bfalse)} \ct \opp{\ap f {\merid(\btrue)}}\\
    &= \lloop \ct \refl{\base}
  \end{align*}
  so this follows easily.
\end{proof}
